\section{统一诸部之后，扩张并非一条直线}

1206年铁木真获成吉思汗称号，蒙古诸部统一。关于称号、诸王关系与早期军政，《元朝秘史》保留了与《元史》不同的叙事层次；二者都不能直接覆盖各部的全部选择。1209年蒙古攻西夏，1211年大举攻金，1215年取中都，说明新的政治共同体很快把草原动员、城战和既有行政资源接在一起。

1218年西辽覆亡，1219年蒙古对花剌子模发动大规模战争，征服范围由东亚延及中亚网络。汉文正史的纪年需要同蒙古、波斯文材料相互校验；地名、译名、军数和因果叙述常有不同，不能把任何一部王朝史当作整个欧亚世界的现场记录。1227年西夏覆亡，同年成吉思汗去世，既是征服节点，也是继承问题的开端。

1229年窝阔台即位，1231年三峰山等战事重创金军主力，1234年蒙古与南宋合攻灭金。短暂联盟并未消除对河南、淮北和道路的竞争。1235年贵族会议决定多方向扩张，显示军事行动需要诸王、将领与后勤网络共同认可；不能把“会议决定”误读为每一条战线均有同样的资源和意志。

\section{大汗继承与汉地征发}

1241年窝阔台去世后，帝国进入摄政与继承协商；1246年贵由即位，1251年蒙哥即位并整肃对手。继承过程既不只是家族私事，也关乎各路线军队、财政收益与任官网络。《元朝秘史》具有口述传统和成书层次，《元史》又由明初编成，两者给出的政治动机都应与文书、碑刻和区域叙事相互限制。

1253年蒙古征服大理，1256年加强汉地户籍、赋税和征发整顿，显示征服后的关键是让人口、土地、工匠和道路能够被记录和调配。制度条文和户籍数字首先说明国家想要什么、如何分类，不自动证明地方已经依同一方式执行。1259年蒙哥在攻宋期间去世，1259--1264年的继承冲突使南方战场与草原政治重新相连。

1260年忽必烈即位，与阿里不哥争位；1261年李璮在山东反叛后被镇压，1264年阿里不哥失败。山东事件说明对汉地官员、军队和财赋的整合并不稳固；继承战争的结束也不等于所有蒙古政治中心从此一致。以“大一统”回望这些年，会抹去多中心帝国中持续存在的协商与竞争。

\section{从大都到江南接收}

1267年忽必烈营建大都，1268年围襄阳、樊城，1271年定国号元，1273年襄樊陷落。这些事件把都城规划、江汉战场和国号语言放在同一时期，却不应当互为简单因果。城市营建需要工匠、物料和运输；围城需要水陆补给、降附者和技术；国号则是对内外政治秩序的声明。

1274年元军分路南下，1276年进入临安，1277年在江南设行省、路府与财赋接收安排，1279年崖山后南宋政权终结。行省设置及其名目可从《元史》与后来的条格材料追索；实际征收、官员流动和地方抵抗要由江南文书、碑刻和区域材料补证。1281年第二次对日远征舰队损失严重，1287年乃颜反叛被平，1294年忽必烈去世、铁穆耳即位；这些节点共同说明，元朝的形成并没有让东北、海域和皇族网络自动服从一个无摩擦的中枢。
